% Nháp Mục 3 (Method) của bài TAIMMA 2027, mẫu LNCS. Không chứa số liệu thực nghiệm.
% Mọi giá trị mặc định khớp mpqda/config.py và docs/taimma/PREREG.md Mục 3; sửa ở đây thì sửa cả hai nơi.
% Cần trong preamble: \usepackage{amsmath,amssymb} \usepackage{algorithm} \usepackage[noend]{algpseudocode}

\section{Method}\label{sec:method}

\subsection{Problem}\label{sec:problem}

Let $X\in\mathbb{R}^{n\times p_0}$ and $y\in\{0,1\}^n$ be the training part of one outer fold after a label-free variance
prefilter that keeps the $p_0=5{,}000$ genes with the largest variance on these samples. Given a budget $k$, the task is to
find a panel $S\subset\{1,\dots,p_0\}$ with $|S|=k$ that maximises a fitness $f(S)$ computed on $(X,y)$ only. The outer test
part is never read during the search. A wrapper is given a fixed number of fitness evaluations; $N_{\mathrm{eval}}$ denotes
the number of distinct panels evaluated, since fitness values are cached by panel.

\paragraph{Fitness.} The training part is split once into three stratified inner folds, which stay fixed during the run and
are shared by all wrappers on the same outer fold. For each inner fold, the genes of $S$ are standardised on the inner-training
samples and an $L_2$-regularised logistic regression ($C=1$, balanced class weights) is fitted. The fitness is
\begin{equation}\label{eq:fitness}
  f(S) = -\frac{1}{3}\sum_{q=1}^{3}\mathrm{LL}_q(S),
\end{equation}
where $\mathrm{LL}_q$ is the log-loss on the validation part of inner fold $q$, with probabilities clipped to
$[10^{-6},1-10^{-6}]$. Log-loss is a strictly proper scoring rule~\cite{gneiting2007}: unlike the area under the ROC curve
(AUROC), it keeps rewarding better-calibrated separation after the two classes are ranked perfectly. With the AUROC fitness,
most wrapper runs reach the ceiling of~1 (Section~\ref{sec:intro}), so the returned panel is one of many tied optima. The AUROC
fitness is kept as a secondary setting to measure this saturation (Section~\ref{sec:results-saturation}).

\subsection{Representation}\label{sec:representation}

An individual is a key vector $x\in(0,1)^{p_0}$~\cite{bean1994}. Its panel is
\begin{equation}\label{eq:decode}
  S(x)=\mathrm{top}_k(x),
\end{equation}
the indices of the $k$ largest keys, ties going to the smaller index. Every individual therefore decodes to a feasible panel
with exactly $k$ genes, and no size penalty is needed.

\paragraph{Band re-encoding.} The fitness depends only on which keys are among the $k$ largest; the other $p_0-k$ keys are not
under selection pressure and do not converge. The scaled difference $F(x_{r_1}-x_{r_2})$ then has magnitude of order $F$ in
every coordinate, far larger than the gaps between the top keys (order $1/p_0$), so a mutant would retain almost nothing of the
panel of its base vector. Every vector that enters the population (initial individuals and accepted trials) is therefore
re-encoded:
\begin{equation}\label{eq:band}
  \mathrm{band}(x)_j=
  \begin{cases}
    \tfrac12+\tfrac12\,\frac{r_j+1/2}{k}, & j\in S(x),\\[2pt]
    \tfrac12\,\frac{r_j+1/2}{p_0-k}, & j\notin S(x),
  \end{cases}
\end{equation}
where $r_j$ is the stable rank of $x_j$ within its group. The panel and the order of keys within each group are unchanged. The
perturbation needed to remove a gene from the panel grows with its rank within the upper band, so $F$ controls how much of the
base panel a mutant keeps. Re-encoding random keys without changing the decoded solution has been used to maintain diversity in a
permutation model-based evolutionary algorithm~\cite{bosman2016}; here the two bands serve a different purpose, namely to let the
DE mutation transmit the panel of its base vector under top-$k$ decoding.

\paragraph{Bound handling.} A mutant coordinate outside $[0,1]$ is moved half-way towards the bound from the target $x_i$:
$v_j>1$ becomes $(1+x_{ij})/2$ and $v_j<0$ becomes $x_{ij}/2$, as in JADE and SHADE~\cite{zhang2009,tanabe2013}. Clipping to
$[0,1]$ would create many keys equal to~1, and the tie rule of~\eqref{eq:decode} would then favour genes with small indices.

\subsection{MPQDA-DE}\label{sec:mpqda}

MPQDA-DE is a random-key differential evolution (DE) with three components: chaotic initialisation, an archive from which the
leader of each mutation is drawn, and a per-individual Q-learning controller of the scale factor. Its output is a consensus
panel computed from the archive. Algorithm~\ref{alg:mpqda} gives the full procedure; terminology follows~\cite{eiben2015}.

\paragraph{Initialisation.} For each individual $i$, $z_i(0)\sim U(0.001,0.999)$ and the keys are successive iterates of the
logistic map, $x_{ij}=z_i(j)$ with $z_i(t+1)=\mu\,z_i(t)\,(1-z_i(t))$ and $\mu=3.99$.

\paragraph{Variation and selection.} For each target $x_i$, a leader $x_{L_i}$ is drawn from the archive and two distinct
individuals $r_1\neq r_2\neq i$ are drawn from the population:
\begin{equation}\label{eq:mutation}
  v_i=\mathrm{repair}\big(x_{L_i}+F_i'\,(x_{r_1}-x_{r_2}),\,x_i\big).
\end{equation}
Binomial crossover with $CR=0.9$ and one forced coordinate $j_{\mathrm{rand}}$ gives the trial $u_i$. The trial replaces the
target if $f(S(u_i))\ge f(S(x_i))$; accepting ties lets the population move across flat regions of the fitness.

\paragraph{Archive and leader.} The archive $A$ holds at most $A_{\max}=15$ panels with their fitness and re-encoded keys. Panels
are compared by the Jaccard distance $d(S_a,S_b)=1-|S_a\cap S_b|/|S_a\cup S_b|$. A candidate identical to a member is ignored.
If a member lies closer than $\delta=0.2$, the nearest such member is replaced only if the candidate has strictly higher fitness;
otherwise the candidate is rejected. Otherwise the candidate is added, and if $|A|>A_{\max}$ the member with the smallest
$0.5\,\widetilde{\mathrm{nov}}+0.5\,\tilde f$ is removed, never the best member. Here $\mathrm{nov}(a)$ is the mean distance of
$a$ to its three nearest members, and the tilde denotes min--max normalisation over the archive. The leader of each individual
is drawn by roulette with weights $0.7\,\widetilde{\mathrm{nov}}+0.3\,\tilde f$, so that leaders come from good panels that lie
in less crowded parts of the archive.

\paragraph{Control of the scale factor.} Each individual keeps its own $F_i$, initialised from $U(0.3,0.8)$ and bounded to
$[0.1,1.0]$, and a state $s_i\in\{0,1\}$: $s_i=0$ if its last trial improved on its target and $s_i=1$ otherwise (initially~1).
A Q-table $Q\in\mathbb{R}^{2\times3}$, shared by all individuals, scores the actions $a\in\{-0.05,0,+0.05\}$. Each generation,
individual $i$ draws $a_i$ from $\mathrm{softmax}(Q[s_i,\cdot]/\tau)$ with $\tau=1$ and uses $F_i'=\mathrm{clip}(F_i+a_i)$. After
the generation is evaluated, with reward $r_i=\mathbb{1}[f(u_i)>f(x_i)]$ and next state $s_i'=1-r_i$,
\begin{equation}\label{eq:q}
  Q[s_i,a_i]\leftarrow Q[s_i,a_i]+\eta\,\big(r_i+\gamma\max_{a}Q[s_i',a]-Q[s_i,a_i]\big),
\end{equation}
with $\eta=0.1$ and $\gamma=0.9$. The new $F_i'$ is kept only if the trial was accepted.

\paragraph{Consensus output.} At termination, members are ranked by fitness and weighted $w_a=|A|-\mathrm{rank}(a)+1$. Each gene
receives the weighted frequency
\begin{equation}\label{eq:consensus}
  \pi_j=\frac{\sum_{a\in A}w_a\,\mathbb{1}[j\in S_a]}{\sum_{a\in A}w_a},
\end{equation}
and the consensus panel $S_c$ holds the $k$ genes with the largest $\pi_j$ (ties: genes of the best member first, then larger
mean key, then smaller index). $S_c$ is evaluated once, so $N_{\mathrm{eval}}\le P(T+1)+1$. Aggregating several near-optimal
panels follows ensemble feature selection~\cite{abeel2010,meinshausen2010} and targets the instability of wrapper panels; the
best individual $S_{\mathrm{best}}$ of the same run is reported as ablation A4.

\begin{algorithm}[t]
\caption{MPQDA-DE for one outer fold and one budget $k$}\label{alg:mpqda}
\begin{algorithmic}[1]
\Require training part $(X,y)$ after the prefilter, budget $k$, seed
\State build the fitness $f$ of~\eqref{eq:fitness} with three fixed inner folds; cache by panel
\State $x_i\gets$ logistic-map keys; draw $F_i\sim U(0.3,0.8)$; $s_i\gets1$; $Q\gets0$ \Comment{$i=1,\dots,P$}
\State $S_i\gets\mathrm{top}_k(x_i)$; $x_i\gets\mathrm{band}(x_i)$; evaluate $f(S_i)$; offer $(S_i,f_i,x_i)$ to $A$
\For{$t=1,\dots,T$}
  \State for each $i$: $a_i\sim\mathrm{softmax}(Q[s_i,\cdot]/\tau)$; $F_i'\gets\mathrm{clip}(F_i+a_i)$
  \State for each $i$: draw leader $L_i$ from $A$ by roulette; draw $r_1\neq r_2\neq i$
  \State for each $i$: $v_i$ by~\eqref{eq:mutation}; $u_i\gets$ binomial crossover of $x_i$ and $v_i$
  \State evaluate $f(\mathrm{top}_k(u_i))$ for all $i$ in one batch
  \State for each $i$: update $Q$ by~\eqref{eq:q}; $s_i\gets s_i'$
  \State for each $i$ with $f(u_i)\ge f(x_i)$: $x_i\gets\mathrm{band}(u_i)$; $F_i\gets F_i'$; offer $(S(u_i),f(u_i),x_i)$ to $A$
\EndFor
\State $S_c\gets$ consensus of $A$ by~\eqref{eq:consensus}; evaluate $f(S_c)$
\State \Return $S_c$ (primary output) and $S_{\mathrm{best}}$
\end{algorithmic}
\end{algorithm}

\paragraph{Implementation.} All random numbers of a generation are drawn before its trials are evaluated, so the trials can be
evaluated in parallel by worker processes (master--slave model~\cite{alba2002}) with results identical to a serial run. This is
an implementation detail and not a contribution. With $P=30$ and $T=50$ the cost is dominated by at most 1{,}531 fitness
evaluations, each fitting three logistic regressions on $k$ genes. Each leader draw and each offer to a full archive recomputes
the novelty at cost $O(A_{\max}^2k)$, so the archive adds $O(P\,A_{\max}^2k)$ operations per generation and the controller
$O(P)$.

% TODO khi có kết quả: thêm câu ngắn về thời gian thực đo ở T6 (Mục 5.6), kèm % source.
% Khoá BibTeX cần có trong Ref.bib: gneiting2007, bean1994, zhang2009, tanabe2013, eiben2015, abeel2010, meinshausen2010, alba2002,
% bosman2016 (doi:10.1145/2908812.2908917) (DOI ở Bảng 11 và Mục 8 của Review-2026-11-14).
% Sửa ngày 14/11 theo T-Y6 (Review-2026-11-14b): bỏ "we" ở dòng 12 và 39; câu về vector hiệu; câu về khoảng cách giữa hai dải;
% chi phí archive O(P A_max^2 k) mỗi thế hệ.
